\chapter{施图姆-刘维尔理论}
\intro{本章建立施图姆-刘维尔理论，为分离变量法提供统一的数学框架。}

\begin{mydefinition}{施图姆-刘维尔问题}
\[
\frac{d}{dx}\left[p(x)\frac{dy}{dx}\right]+[\lambda w(x)-q(x)]y=0
\]
配合适当的边界条件。
\end{mydefinition}

\begin{theorem}[特征值性质]
\begin{enumerate}
	\item 特征值为实数且可排序：$\lambda_1<\lambda_2<\cdots$。
	\item 对应不同特征值的特征函数关于权函数 $w(x)$ 正交：
	\[
	\int_a^b\phi_m(x)\phi_n(x)w(x)\,dx=0,\quad m\neq n
	\]
\end{enumerate}
\end{theorem}

\begin{theorem}[完备性]
特征函数系 $\{\phi_n\}$ 构成 $L^2_w[a,b]$ 空间的完备正交基。任意 $f\in L^2_w$ 可展开为广义傅里叶级数。
\end{theorem}

\begin{remark}
Sturm-Liouville 理论将分离变量法的结果统一：热传导、波动、拉普拉斯方程的特征函数展开都是 S-L 理论的特例。
\end{remark}
